\section{案例研究：Vision-guided Autonomy}
\subsection{现实场景中的任务拆解}
Slide 09 展示了一个 vision-guided manipulation case，包含 grasp、transport、place 三个模块，每个模块又嵌套自主 reward、reset 与 goal manager。讲者强调：真实场景常常有 long-tailed failure mode，需要在 pipeline 中引入 anomaly detection。

\begin{figure}[H]
\centering
\includegraphics[width=0.85\textwidth]{slides-images/slide-09.jpg}
\caption{Slide 9 展示 vision-guided manipulation 的工作流，强调各 module 如何协同。}
\end{figure}

\begin{knowledgebox}{Vision-guided Autonomy 切分}
\begin{itemize}
\item Grasp module：使用 visual affordance 预测 grasp point；\\
\item Transport module：通过 goal-conditioned policy 将物体移动；\\
\item Place module：利用 success classifier 判断目标放置是否完成。
\end{itemize}
\end{knowledgebox}

\subsection{可观察性与异常响应}
讲者建议在每个阶段记录 anomaly metric：vision 模型 confidence drop、reset policy failure rate、goal coverage drop。把这些度量填入 dashboard，可以 trigger automated rollback 或 reward recalibration。

\begin{importantbox}{异常响应矩阵}
\begin{tabularx}{\textwidth}{lX}
\toprule
异常类型 & 响应 \\
\midrule
Vision confidence drop & 进入 safe mode，降低 exploration rate； \\
Reset failure & 立即触发 manual inspection ，并回滚到最近 stable checkpoint； \\
Goal coverage shrink & 扩展 goal replay buffer，重新 thaw failed goal cluster； \\
\bottomrule
\end{tabularx}
\end{importantbox}

\subsection{本章小结}
通过 vision-guided manipulation 案例，展示 Autonomy pipeline 在线下与线上部署中的可观察性与异常响应策略。
\newpage

